% Mechanistic Validity — NEMI 2026 poster, Option A
% "What Five Fields Already Knew"
% 3-column, 3-row. Text and tables only (no figures).
%
%   cd poster && lualatex mechval_poster_option_a.tex

\documentclass[25pt, a0paper, portrait]{tikzposter}

\usepackage{amsmath,amssymb}
\usepackage{array}
\usepackage{booktabs}
\usepackage{graphicx}
\usepackage{enumitem}
\usepackage{microtype}

% ── Color palette ───────────────────────────────────────
\definecolor{confirmed}{HTML}{27AE60}
\definecolor{partial}{HTML}{2980B9}
\definecolor{inconclusive}{HTML}{F39C12}
\definecolor{untested}{HTML}{95A5A6}
\definecolor{disconfirmed}{HTML}{E74C3C}
\definecolor{darktext}{HTML}{2C3E50}
\definecolor{blockbg}{HTML}{FFFFFF}
\definecolor{tierProposed}{HTML}{95A5A6}
\definecolor{tierCausal}{HTML}{F39C12}
\definecolor{tierSupported}{HTML}{2980B9}
\definecolor{tierTriangulated}{HTML}{27AE60}
\definecolor{tierValidated}{HTML}{8E44AD}
\definecolor{tierDisconfirmed}{HTML}{E74C3C}
\definecolor{novelgreen}{HTML}{1ABC9C}
\definecolor{fieldblue}{HTML}{34495E}

% ── Theme ───────────────────────────────────────────────
\usetheme{Default}

\definecolorstyle{MechValColors}{
  \definecolor{colorOne}{HTML}{2C3E50}
  \definecolor{colorTwo}{HTML}{3498DB}
  \definecolor{colorThree}{HTML}{ECF0F1}
}{
  \colorlet{backgroundcolor}{colorThree}
  \colorlet{framecolor}{colorOne}
  \colorlet{titlefgcolor}{white}
  \colorlet{titlebgcolor}{colorOne}
  \colorlet{blocktitlebgcolor}{colorTwo}
  \colorlet{blocktitlefgcolor}{white}
  \colorlet{blockbodybgcolor}{blockbg}
  \colorlet{blockbodyfgcolor}{darktext}
  \colorlet{notefgcolor}{darktext}
  \colorlet{notebgcolor}{inconclusive!15}
  \colorlet{notefrcolor}{inconclusive}
}
\usecolorstyle{MechValColors}

\defineblockstyle{MechValBlock}{
  titlewidthscale=1.0, bodywidthscale=1.0, titleleft,
  titleoffsetx=0pt, titleoffsety=0pt,
  bodyoffsetx=0pt, bodyoffsety=0pt, bodyverticalshift=0pt,
  roundedcorners=12, linewidth=2pt, titleinnersep=10pt, bodyinnersep=15pt
}{
  \draw[rounded corners=\blockroundedcorners, color=blocktitlebgcolor,
        fill=blockbodybgcolor, line width=\blocklinewidth]
    (blockbody.south west) rectangle (blocktitle.north east);
  \ifBlockHasTitle
    \draw[rounded corners=\blockroundedcorners,
          fill=blocktitlebgcolor, color=blocktitlebgcolor,
          line width=\blocklinewidth]
      (blocktitle.south west) rectangle (blocktitle.north east);
  \fi
}
\useblockstyle{MechValBlock}

\definetitlestyle{MechValTitle}{
  width=\paperwidth, roundedcorners=0, linewidth=0pt, innersep=20pt,
  titletotopverticalspace=15mm, titletoblockverticalspace=25mm
}{
  \draw[fill=titlebgcolor, color=titlebgcolor]
    (\titleposleft, \titleposbottom)
    rectangle (\titleposright, \titlepostop);
}
\usetitlestyle{MechValTitle}

\title{%
  \parbox{0.92\linewidth}{\centering\bfseries
    Mechanistic Validity\\[6pt]
    What Five Fields Already Knew About Validating Mechanistic Claims}}
\author{Elliot Tower\quad\texttt{elliot@elliottower.ai}}
\institute{}

\begin{document}

\maketitle[width=\paperwidth]

% ════════════════════════════════════════════════════════
% ROW 1 — BORROWED + NOVEL + MINIMUM RULE
% ════════════════════════════════════════════════════════
\begin{columns}
\column{0.33}

{
\colorlet{blocktitlebgcolor}{fieldblue}
\block{Borrowed from Five Fields}{
  Every field that makes mechanistic claims has a validity
  framework. We imported the criteria that apply to neural
  network mechanisms.

  \vspace{5mm}
  \begin{center}
  \renewcommand{\arraystretch}{1.25}
  {\small
  \begin{tabular}{@{}>{\raggedright\arraybackslash}p{2.8cm}
                    >{\raggedright\arraybackslash}p{3.8cm}l@{}}
  \toprule
  \textbf{Field} & \textbf{Contribution} & \textbf{Criteria} \\
  \midrule
  Philosophy of science
    & Falsifiability, severe testing
    & C1, I5 \\
  Measurement theory
    & Construct + discriminant validity
    & C3, C4 \\
  Causal inference
    & Sensitivity analysis, E-values
    & I7, I8 \\
  Neuroscience
    & Constitutive relevance, levels
    & V1, V2 \\
  Genetics
    & Complementation, rescue
    & C6, I10 \\
  \bottomrule
  \end{tabular}}
  \end{center}

  \vspace{4mm}
  The criteria are almost entirely imported. What is new is the
  object being assessed and the treatment of a claim as a
  family of readings at different strengths.
}
}

\column{0.34}

{
\colorlet{blocktitlebgcolor}{novelgreen}
\block{What We Add: Interpretive Validity}{
  Other fields do not need interpretive validity because their
  mechanisms are not \textbf{named by the experimenter}.

  \vspace{5mm}
  A treatment effect cannot be pitched at the wrong level of
  description. A test score either measures its construct or
  does not. A claim about how a neural network works can be
  well-measured, causally established, and still wrong about
  what its components do---because the label asserts more than
  the intervention supports.

  \vspace{5mm}
  \begin{tabular}{@{}lp{0.72\linewidth}@{}}
  \toprule
  \textbf{ID} & \textbf{Criterion} \\
  \midrule
  V1 & Level of description declared \\
  V2 & Evidence matches the declared level \\
  V3 & Alternative level considered \\
  V4 & Labels do not exceed evidence \\
  V5 & Scope explicitly stated \\
  \bottomrule
  \end{tabular}

  \vspace{5mm}
  Also new: the dependency chain (construct validity caps
  everything) and the minimum rule.
}
}

\column{0.33}

{
\colorlet{blocktitlebgcolor}{disconfirmed}
\block{The Minimum Rule}{
  {\Large\bfseries A claim is as strong as its weakest
  validity dimension.}

  \vspace{6mm}
  \renewcommand{\arraystretch}{1.4}
  \begin{tabular}{@{}l>{\raggedright\arraybackslash}p{0.68\linewidth}@{}}
  \textcolor{tierTriangulated}{\Large$\blacksquare$}
    & \textbf{Construct}: Is the circuit a natural kind? \\
  \textcolor{tierSupported}{\Large$\blacksquare$}
    & \textbf{Measurement}: Does the metric measure what
      it claims? \\
  \textcolor{tierCausal}{\Large$\blacksquare$}
    & \textbf{Internal}: Is the causal claim warranted? \\
  \textcolor{inconclusive}{\Large$\blacksquare$}
    & \textbf{External}: Does the result generalize? \\
  \textcolor{novelgreen}{\Large$\blacksquare$}
    & \textbf{Interpretive}: Is the label justified? \\
  \end{tabular}

  \vspace{6mm}
  Not the average. Not the best. \textbf{The weakest.}

  \vspace{4mm}
  No amount of causal evidence rescues a claim whose target
  concept is incoherent. Construct validity is a ceiling on
  all the rest.
}
}

\end{columns}

% ════════════════════════════════════════════════════════
% ROW 2 — SIXTEEN CLAIMS (wide table)
% ════════════════════════════════════════════════════════
\begin{columns}
\column{0.50}

\block{Sixteen Claims from Fifteen Papers}{
  {\small
  \renewcommand{\arraystretch}{1.2}
  \begin{tabular}{@{}>{\raggedright\arraybackslash}p{3.2cm}
                    >{\raggedright\arraybackslash}p{5.5cm}
                    >{\raggedright\arraybackslash}p{4.0cm}@{}}
    \toprule
    \textbf{Verdict} & \textbf{Claim} & \textbf{Scope} \\
    \midrule
    \textcolor{tierTriangulated}{\textbf{Triangulated}}
      & Induction heads: token copying
      & 1L--70B \\
    \midrule
    \textcolor{tierSupported}{\textbf{Mech.\ Supported}}
      & Fourier mechanism (mod.\ addition)
      & 1L ReLU \\
      & Greater-Than circuit
      & GPT-2 Small \\
      & Copy Suppression (L10H7)
      & GPT-2 Small \\
      & Superposition (toy networks)
      & 4 ReLU nets \\
      & Steering vector (refusal)
      & 13 models \\
      & Global Workspace / J-space
      & 4 Claude \\
      & Successor Heads
      & Pythia-1.4B \\
    \bottomrule
  \end{tabular}}
}

\column{0.50}

\block{~}{
  {\small
  \renewcommand{\arraystretch}{1.2}
  \begin{tabular}{@{}>{\raggedright\arraybackslash}p{3.2cm}
                    >{\raggedright\arraybackslash}p{5.5cm}
                    >{\raggedright\arraybackslash}p{4.0cm}@{}}
    \toprule
    \textbf{Verdict} & \textbf{Claim} & \textbf{Scope} \\
    \midrule
    \textcolor{tierCausal}{\textbf{Causally Suggestive}}
      & Othello World Model
      & Othello-GPT \\
      & Docstring Circuit
      & 4L attn-only \\
      & Gender Bias Circuits
      & GPT-2 + 5 \\
      & IOI Circuit
      & GPT-2 Small \\
    \midrule
    \textcolor{tierProposed}{\textbf{Proposed}}
      & Probing Classifiers
      & review \\
    \midrule
    \textcolor{tierDisconfirmed}{\textbf{Disconfirmed}}
      & Knowledge Neurons
      & BERT-base \\
      & SAE features as canonical units
      & Pythia \\
      & Induction heads: general ICL
      & $>$1B \\
    \bottomrule
  \end{tabular}}

  \vspace{4mm}
  \textbf{No claim reaches Validated.} The highest tier achieved
  is Triangulated, by one claim---and only the narrow reading.
}

\end{columns}

% ════════════════════════════════════════════════════════
% ROW 3 — SAME PAPER TWO VERDICTS + CAPS + TAKEAWAY
% ════════════════════════════════════════════════════════
\begin{columns}
\column{0.33}

{
\colorlet{blocktitlebgcolor}{tierTriangulated}
\block{Same Paper, Two Verdicts}{
  Olsson et al.\ (2022) make two claims. The framework
  separates them.

  \vspace{5mm}
  \textbf{\textcolor{tierTriangulated}{Token copying:
  Triangulated}}

  14 Confirmed, 15 Partially confirmed, 6 Untested.
  Replicates to 70B: removing the top 10\% of induction heads
  collapses 8-shot accuracy from 80.6\% to 2.5\%. Double
  dissociation obtained by Feucht et al.

  \vspace{5mm}
  \textbf{\textcolor{tierDisconfirmed}{General ICL:
  Disconfirmed}}

  6 Confirmed, 14 Partially confirmed, 2 Inconclusive,
  8 Untested, \textbf{5 Disconfirmed}. At 70B the heads
  carrying abstract reasoning are disjoint from induction
  heads ($r = 0.11$). The authors anticipated this in
  Argument~6.

  \vspace{5mm}
  One paper. Same evidence base. Two tiers.
}
}

\column{0.34}

{
\colorlet{blocktitlebgcolor}{tierCausal}
\block{What Caps Each Verdict}{
  Every capped verdict names the criterion that caps it.

  \vspace{5mm}
  \textbf{IOI Circuit} --- Faithfulness is method-conditional.
  The headline 87\% spans below 0\% to over 100\% across six
  methodological choices (Miller et al.).
  Rival head sets reproduce the behavior at 100\%
  accuracy with 4.1\% edge overlap.

  \vspace{5mm}
  \textbf{SAE Features} --- Seed-dependent: cross-seed
  agreement falls from 94\% at 3k latents to 30\% at 131k.
  Interpretability scores do not distinguish trained
  transformers from random ones. Features localize causal
  variables no better than raw neurons.

  \vspace{5mm}
  \textbf{Othello ``World Model''} --- The label exceeds
  the evidence. Nonlinear probes recover board state at 1.7\%
  error; linear probes at 20.4\%. The label ``world model''
  asserts more than ``linearly decodable.'' Criterion V4.

  \vspace{5mm}
  \textbf{Gender Bias} --- Necessity is untestable: every
  estimand is a substitution, never a removal. No capability
  measure is defined, so specificity has no second quantity.
}
}

\column{0.33}

{
\colorlet{blocktitlebgcolor}{colorOne}
\block{Takeaway}{
  {\Large\bfseries Report the tier your evidence supports,
  and name the criterion that caps it.}

  \vspace{6mm}
  \begin{itemize}[leftmargin=*, itemsep=4mm]
    \item \textbf{36} criteria across \textbf{5} validity types,
      in dependency order
    \item \textbf{16} claims audited from \textbf{15} papers
    \item \textbf{0} claims reach Validated
    \item One paper appears at both ends: Triangulated and
      Disconfirmed
    \item The criteria are imported; what is new is the object
      (a mechanistic claim) and interpretive validity
  \end{itemize}

  \vspace{8mm}
  \begin{center}
  \texttt{elliot@elliottower.ai}

  \vspace{3mm}
  {\small Framework, criteria, and all sixteen evidence cards:}

  \texttt{mechanisticresearch.org}
  \end{center}
}
}

\end{columns}

\end{document}
